% !TeX root = ../main.tex

\begin{denotation}[4.0cm]

  %% ──────────────────────────────────────────────
  %% 一、模型、统计检验与通用缩略语
  %% ──────────────────────────────────────────────
  \item[MSAS-GNN] 多尺度自适应稀疏图神经网络（Multi-Scale Adaptive Sparse Graph Neural Network，本文提出）
  \item[SDGNN] 稀疏分解图神经网络（Sparse Decomposition of Graph Neural Networks，Hu等，2024）
  \item[GNN] 图神经网络（Graph Neural Network）
  \item[GCN] 图卷积网络（Graph Convolutional Network）
  \item[GAT] 图注意力网络（Graph Attention Network）
  \item[GIN] 图同构网络（Graph Isomorphism Network）
  \item[SGC] 简单图卷积（Simple Graph Convolution）
  \item[MPNN] 消息传递神经网络（Message Passing Neural Network）
  \item[APPNP] 近似个性化 PageRank 传播图网络（Approximate Personalized PageRank Propagation of Neural Predictions）
  \item[SIGN] 可扩展 Inception 图神经网络（Scalable Inception Graph Neural Network）
  \item[FAGCN] 频率自适应图卷积网络（Frequency Adaptive Graph Convolution Network）
  \item[GPR-GNN] 广义 PageRank 图神经网络（Generalized PageRank Graph Neural Network）
  \item[H2GCN] 面向异配图的图卷积网络（Beyond Homophily Graph Convolutional Network）
  \item[Geom-GCN] 几何聚合图卷积网络（Geometric Graph Convolutional Network）
  \item[GLNN] 无图神经网络（Graph-less Neural Network）
  \item[GloGNN] 全局同配建模图神经网络（Global Homophily based Graph Neural Network）
  \item[LINKX] 面向大规模异配图的线性图学习方法（Linear Graph Model for Heterophilous Graphs）
  \item[FastGCN] 快速采样图卷积网络（Fast Graph Convolutional Network）
  \item[GraphSAGE] 归纳式大规模图表示学习方法（Graph Sample and Aggregate）
  \item[GraphSAINT] 基于图采样的归纳学习方法（Graph Sampling Based Inductive Learning Method）
  \item[DropEdge] 随机边丢弃正则化方法（DropEdge）
  \item[PPRGo] 基于近似 Personalized PageRank 的可扩展图神经网络（Scaling Graph Neural Networks with Approximate PageRank）
  \item[NodeFormer] 可扩展图结构学习 Transformer（Scalable Graph Structure Learning Transformer）
  \item[DIFFormer] 基于能量约束扩散的可扩展图 Transformer（Scalable Graph Transformers Induced by Energy Constrained Diffusion）
  \item[SGFormer] 简化图 Transformer（Simplified Graph Transformer）
  \item[NAGphormer] 邻域聚合图 Transformer（Neighborhood Aggregation Graph Transformer）
  \item[LARS] 最小角回归（Least Angle Regression）
  \item[MLP] 多层感知机（Multilayer Perceptron）
  \item[GSP] 图信号处理（Graph Signal Processing）
  \item[OGB] 开放图基准（Open Graph Benchmark）
  \item[Wilcoxon] Wilcoxon 符号秩检验（Wilcoxon Signed-Rank Test）
  \item[GPU] 图形处理器（Graphics Processing Unit）
  \item[Adam] 自适应矩估计优化器（Adaptive Moment Estimation）
  \item[$t$-SNE] $t$ 分布随机近邻嵌入（$t$-distributed Stochastic Neighbor Embedding）
  
  %% ──────────────────────────────────────────────
  %% 二、图结构与基础任务符号
  %% ──────────────────────────────────────────────
  \item[$G=(V,E,X)$] 属性图，$V$ 为节点集，$E$ 为边集，$X$ 为节点特征矩阵
  \item[$n$] 节点总数，$n=|V|$
  \item[$m$] 邻接矩阵非零元数，统一记为 $m:=\mathrm{nnz}(A)=2|E|$
  \item[$D$] 节点输入特征维度
  \item[$d$] 节点嵌入维度或隐藏维度
  \item[$C$] 下游任务类别数
  \item[$A$] 图邻接矩阵
  \item[$D_{\mathrm{deg}}$] 度矩阵，对角元素为节点度数
  \item[$d_i$] 节点 $i$ 的度
  \item[$\bar{d}$] 全图平均度，$\bar{d}=m/n$
  \item[$V_{\mathrm{train}}$] 训练节点集合，满足 $V_{\mathrm{train}}\subseteq V$
  \item[$f_{\mathrm{cls}}(\cdot;W_{\mathrm{cls}})$] 节点分类头或下游分类映射
  \item[$W_{\mathrm{cls}}$] 分类头参数
  \item[$\hat{y}_i$] 节点 $i$ 的预测类别分布，$\hat{y}_i\in\Delta^{C-1}$
  \item[$\Delta^{C-1}$] $C$ 维概率单纯形
  \item[$\mathcal{N}(i)$] 节点 $i$ 的一跳邻域
  \item[$\mathcal{N}_l(i)$] 节点 $i$ 的 $l$ 跳邻域
  \item[$\mathcal{N}_L(i)$] 节点 $i$ 的 $L$ 跳感受野邻域
  \item[$\mathcal{R}_l(i)$] 节点 $i$ 的第 $l$ 跳邻居环层，$\mathcal{R}_l(i)=\mathcal{N}_l(i)\setminus \mathcal{N}_{l-1}(i)$
  \item[$d_G(i,j)$] 图 $G$ 上节点 $i$ 与 $j$ 的最短路径距离

  %% ──────────────────────────────────────────────
  %% 三、谱图理论与离线预处理符号
  %% ──────────────────────────────────────────────
  \item[$\tilde{\mathcal{L}}$] 归一化图拉普拉斯矩阵，$\tilde{\mathcal{L}}=I-D_{\mathrm{deg}}^{-1/2}AD_{\mathrm{deg}}^{-1/2}$
  \item[$U,\Lambda$] $\tilde{\mathcal{L}}$ 的谱分解矩阵，$\tilde{\mathcal{L}}=U\Lambda U^\top$
  \item[$\lambda_k$] 归一化拉普拉斯的第 $k$ 个特征值，满足 $0\le \lambda_k\le 2$
  \item[$\lambda_{\mathrm{gap}}$] 谱间隙，即最小非零特征值；在连通图上有 $\lambda_{\mathrm{gap}}=\lambda_1$
  \item[$\hat{\lambda}_{\mathrm{gap}}$] 由 Lanczos 预处理得到的谱间隙近似值
  \item[$K_{\mathrm{eig}}$] Lanczos 迭代中保留的特征对数量
  \item[$K_{\mathrm{pair}}$] 预处理中实际返回并用于后续计算的特征对数量
  \item[$K_{\mathrm{eff}}$] Lanczos 迭代的有效终止步数
  \item[$\tau_{\mathrm{tol}}$] Lanczos 迭代提前停止所用数值容差
  \item[$T_{K_{\mathrm{eff}}}$] Lanczos 过程生成的三对角矩阵
  \item[$g(\tilde{\mathcal{L}})$] 作用于归一化拉普拉斯的谱滤波器
  \item[$K_{\mathrm{poly}}$] 多项式谱滤波器的阶数
  \item[$\bar{\mathcal{L}}$] 平移后的拉普拉斯矩阵，$\bar{\mathcal{L}}:=\tilde{\mathcal{L}}-I$
  \item[$a_r$] 多项式谱滤波器单项式展开的第 $r$ 阶系数
  \item[$S$] 无自环对称归一化传播算子，$S=D_{\mathrm{deg}}^{-1/2}AD_{\mathrm{deg}}^{-1/2}$
  \item[$S_{\mathrm{lazy}}$] 懒化传播算子，$S_{\mathrm{lazy}}=(I+S)/2$
  \item[$P$] 分层信息衰减分析中的线性传播算子，可取 $S$ 或 $S_{\mathrm{lazy}}$
  \item[$\mu_k(P)$] 传播算子 $P$ 的第 $k$ 个特征值
  \item[$\mu_\star(P)$] 传播算子 $P$ 在 $\mathbf{1}^{\perp}$ 子空间上的第二谱半径
  \item[$h_i^{[l]}$] 节点 $i$ 在第 $l$ 次线性传播后的累计表示
  \item[$\Delta h_i^{[l]}$] 第 $l$ 次传播相对第 $l-1$ 次传播的边际表示增量
  \item[$e_i$] 第 $i$ 个标准基向量
  \item[$C_0$] 分层信息衰减界中的全局保守常数，$C_0=\|(P-I)X\|_F$
  \item[$\sigma$] 原图与稀疏图之间的谱相似性参数，$\sigma\ge 1$
  \item[$\widetilde{\sigma}_{\mathrm{proxy}}$] 基于 Lanczos 的 $\sigma$-谱相似性代理估计量
  \item[$K_{\mathrm{Lanczos}}$] 估计 $\widetilde{\sigma}_{\mathrm{proxy}}$ 时采用的 Lanczos 固定步数
  \item[$\hat{G}$] 由稀疏结构诱导的稀疏图
  \item[$\tilde{\mathcal{L}}_G$] 原图 $G$ 的归一化拉普拉斯矩阵
  \item[$\tilde{\mathcal{L}}_{\hat{G}}$] 稀疏图 $\hat{G}$ 的归一化拉普拉斯矩阵
  \item[$\Delta_{\mathcal{L}}$] 稀疏图与原图归一化拉普拉斯之差，$\Delta_{\mathcal{L}}=\tilde{\mathcal{L}}_{\hat{G}}-\tilde{\mathcal{L}}_G$
  \item[$M$] 谱相似性推导中的半正定上界矩阵，$M=(\sigma-1)\tilde{\mathcal{L}}_G$
  \item[$\lambda_{\max}(M)$] 矩阵 $M$ 的最大特征值
  \item[$A\preceq B$] Loewner 半序，表示 $B-A$ 为半正定矩阵
  \item[$\eta$] 原图与稀疏图平移拉普拉斯之差的谱范数，$\eta=\|\bar{\mathcal{L}}_G-\bar{\mathcal{L}}_{\hat{G}}\|_2$

  %% ──────────────────────────────────────────────
  %% 四、图复杂度指标体系
  %% ──────────────────────────────────────────────
  \item[$w_k(i)$] 节点 $i$ 对第 $k$ 个频率分量的权重；正文主线取等权形式
  \item[$E_{\mathrm{spectral}}(i)$] 节点谱能量，衡量节点在各谱分量上的参与强度
  \item[$H(i)$] 节点 $i$ 的局部图熵
  \item[$H_{\mathrm{norm}}(i)$] 节点 $i$ 的归一化局部图熵，$H_{\mathrm{norm}}(i)\in[0,1]$
  \item[$C_{\mathrm{deg}}(i)$] 节点 $i$ 的度中心性
  \item[$\operatorname{core}(i)$] 节点 $i$ 的 $k$-core 指数
  \item[$k_{\max}$] 全图节点 $k$-core 指数的最大值，$k_{\max}=\max_{j\in V}\operatorname{core}(j)$
  \item[$w_{ij}$] 构造局部图熵时节点 $i$ 与 $j$ 之间的边权
  \item[$Z_i$] 节点 $i$ 局部熵计算中的归一化常数，$Z_i=\sum_{j\in \mathcal{N}(i)}w_{ij}$
  \item[$\hat{\sigma}^2$] 高斯核边权中使用的经验带宽估计
  \item[$\gamma_x$] 高斯核带宽参数，满足 $w_{ij}=\exp(-\gamma_x\|x_i-x_j\|_2^2)$
  \item[$\mathrm{CV}(\{k_i\})$] 节点稀疏预算分布的变异系数，用于辅助分析节点间差异性

  %% ──────────────────────────────────────────────
  %% 五、MSAS-GNN 稀疏分解与优化目标
  %% ──────────────────────────────────────────────
  \item[$L$] 图传播层数或最大跳数
  \item[$H^{(l)}$] 第 $l$ 层节点表示矩阵，$H^{(l)}\in\mathbb{R}^{n\times d}$
  \item[$H^*$] 教师 GNN 的目标表示矩阵，$H^*\in\mathbb{R}^{n\times d}$
  \item[$h_i^*$] 教师 GNN 在节点 $i$ 处的目标表示向量
  \item[$\phi(\cdot;W_\phi)$] 节点特征变换函数
  \item[$W_\phi$] 特征变换函数的可学习参数
  \item[$W_\phi^*$] 交替优化初始化时使用的预训练特征变换参数
  \item[$\Phi=\phi(X;W_\phi)$] 特征变换后的节点表示矩阵
  \item[$\tilde{\Phi}$] 行 $\ell_2$ 归一化后的特征变换矩阵
  \item[$\tilde{\Phi}^{\mathrm{cache}}$] 推理期缓存的归一化特征变换矩阵
  \item[$\theta_i$] 节点 $i$ 对应的稀疏权重向量
  \item[$\Theta$] 全图稀疏权重矩阵，$\Theta\in\mathbb{R}^{n\times n}$，其第 $i$ 列为节点 $i$ 的稀疏权重向量 $\theta_i$
  \item[$\Theta^{\mathrm{fixed}}$] 推理期固化后的稀疏权重矩阵
  \item[$\hat{h}(i)$] 节点 $i$ 的稀疏推理表示
  \item[$\hat{H}$] 稀疏分解推理输出矩阵，$\hat{H}=\Theta^\top\tilde{\Phi}$
  \item[$m_{\mathrm{sparse}}$] 稀疏结构的非零元数，$m_{\mathrm{sparse}}:=\mathrm{nnz}(\Theta)$
  \item[$k_i$] 节点 $i$ 最终保留的非零邻居数，常记为 $k_i=\|\theta_i\|_0$
  \item[$\bar{k}$] 全图平均非零权重数或平均稀疏预算，$\bar{k}=\frac{1}{n}\sum_{i=1}^n k_i$；按预算上限估计复杂度时，可用 $\frac{1}{n}\sum_i k_i^{\mathrm{total}}$ 近似
  \item[$\lambda_W$] 特征变换参数的 $\ell_2$ 正则系数
  \item[$\lambda_{\mathrm{reg}}$] SDGNN 使用的全局统一正则系数
  \item[$\mathcal{L}_{\mathrm{task}}$] 下游任务损失，在节点分类实验中取交叉熵损失
  \item[$\mathcal{E}_{\mathrm{approx}}$] 表示近似误差，定义为 $\mathcal{E}_{\mathrm{approx}}=\frac{1}{\sqrt{n}}\|H^*-\hat{H}\|_F$
  \item[$\varepsilon_i^{\mathrm{rep}}$] 节点 $i$ 的表示误差范数，$\varepsilon_i^{\mathrm{rep}}=\|h_i^*-\hat{h}_i\|_2$
  \item[$\mathcal{S}(t)$] 在误差阈值 $t$ 下的节点子集，用于第3章的设计准则分析
  \item[$\varepsilon_0$] 第3章性能保持准则中的总体表示误差容许阈值

  %% ──────────────────────────────────────────────
  %% 六、节点级自适应正则与跳距预算
  %% ──────────────────────────────────────────────
  \item[$\tau(i)$] 节点级正则系数，取值区间为 $[\tau_{\min},\tau_{\mathrm{base}}]$
  \item[$\tau_{\mathrm{base}}$] 节点级正则系数的基准上界
  \item[$\tau_{\min}$] 节点级正则系数的下界
  \item[$\tau_{\mathrm{mono}}(i)$] 仅由节点谱能量驱动的单因子节点级正则系数
  \item[$\beta_{\tau}$] 节点谱能量到正则系数映射中的响应强度系数
  \item[$\gamma$] 多因子节点级正则系数中度中心性项的权重
  \item[$\delta$] 多因子节点级正则系数中 $k$-core 项的权重
  \item[$\omega_H$] 多因子节点级正则系数中归一化局部熵项的权重
  \item[$c_E$] 节点谱能量阈值的尺度系数
  \item[$E_{\mathrm{threshold}}$] 节点谱能量阈值，定义为 $E_{\mathrm{threshold}}=c_E\cdot \operatorname{median}_{j\in V}E_{\mathrm{spectral}}(j)$
  \item[$k_i^{\mathrm{total}}$] 节点 $i$ 的总稀疏预算
  \item[$p_l$] 第 $l$ 跳的保留率，满足 $p_l\in(0,1]$
  \item[$\varrho$] 跳距保留率的几何衰减公比
  \item[$p_{\mathrm{base}}$] 谱间隙自适应跳距分配中的第1跳基准保留率
  \item[$p_{\min}$] 跳距保留率下界，用于避免高跳距预算被完全压缩
  \item[$\kappa$] 谱间隙自适应跳距保留率中的响应系数，$\kappa>0$
  \item[$\rho$] 跳距预算分配中的全局预算缩放系数，$\rho\in(0,1]$
  \item[$\tilde{k}_i^{(l)}$] 节点 $i$ 第 $l$ 跳环层的连续预算上界
  \item[$\operatorname{RoundQuota}(\cdot)$] 将连续预算上界转化为整数预算配额的舍入算子
  \item[$\xi$] 跳距预算分配中的整数化控制参数
  \item[$k_i^{(l)}$] 节点 $i$ 在第 $l$ 跳邻居环层上的预算配额，满足 $\sum_{l=1}^{L}k_i^{(l)}=k_i^{\mathrm{total}}$
  \item[$\varepsilon_l$] 第 $l$ 跳传播允许的信息损失预算
  \item[$\varepsilon_{\mathrm{total}}$] 第5章条件化误差上界中的整体误差目标

  %% ──────────────────────────────────────────────
  %% 七、算法实现与复杂度分析符号
  %% ──────────────────────────────────────────────
  \item[$C_i$] 节点 $i$ 的候选节点集合，$C_i=\{i\}\cup \bigcup_{l=1}^{L}\mathcal{R}_l(i)$
  \item[$\Omega_i$] 节点 $i$ 的目标表示向量，$\Omega_i:=H_{i,:}^{*}$
  \item[$\tilde{\phi}_i$] 行归一化特征矩阵 $\tilde{\Phi}$ 的第 $i$ 行向量
  \item[$\tilde{\Phi}_{\mathcal{R}_l(i)}$] $\tilde{\Phi}$ 在节点 $i$ 第 $l$ 跳邻居环层上的行子矩阵
  \item[$\theta_i^{(0)}$] 节点 $i$ 的第0跳自特征通道系数
  \item[$\theta_i^{(l)}$] 节点 $i$ 在第 $l$ 跳邻居环层上的稀疏系数向量
  \item[$\theta_{i,\mathrm{share}}^{(l)}$] 共享目标方案下节点 $i$ 第 $l$ 个候选块的稀疏系数
  \item[$r_i^{(l)}$] 节点 $i$ 完成第0跳至第 $l$ 跳拟合后的剩余残差
  \item[$\Theta_{j,i}$] 稀疏权重矩阵 $\Theta$ 的第 $j$ 行第 $i$ 列元素
  \item[$\operatorname{AdamStep}(\cdot)$] 对参数执行一次 Adam 更新的算子
  \item[$T_{\mathrm{epoch}}$] 第5章主循环中的外层交替优化轮数
  \item[$T_W$] 每轮外层迭代中 Phase $W$ 的更新步数
  \item[$\eta_W$] Phase $W$ 中的学习率
  \item[$q$] 早停窗口长度
  \item[$O(Lmd)$] 标准消息传递 GNN 在全图一次前向口径下的主项复杂度
  \item[$O(n\bar{k}d)$] MSAS-GNN / SDGNN 在推理阶段的主项复杂度口径

  %% ──────────────────────────────────────────────
  %% 八、实验评估符号
  %% ──────────────────────────────────────────────
  \item[$t_{\mathrm{dense}}$] 稠密基线模型推理时间
  \item[$t_{\mathrm{sparse}}$] MSAS-GNN 稀疏推理时间
  \item[$\mathrm{Speedup}$] 推理加速比，$\mathrm{Speedup}=t_{\mathrm{dense}}/t_{\mathrm{sparse}}$
  \item[$t_{\mathrm{pre}}$] 总预处理时间，包括指标计算、LARS 求解与结构固化
  \item[$Q_{\mathrm{be}}$] 预处理的 break-even 调用次数，$Q_{\mathrm{be}}=t_{\mathrm{pre}}/(t_{\mathrm{dense}}-t_{\mathrm{sparse}})$
  \item[$B_t$] 第 $t$ 轮 mini-batch 中采样的节点子集
  \item[\texttt{batch\_size}] mini-batch 训练或批次推理中的批大小
  \item[$p$ 值] 统计检验中的显著性概率值
  \item[$r_{\mathrm{Pearson}}$] Pearson 相关系数，用于描述实验可视化中的线性相关程度

\end{denotation}